% Formalización nueva integrada en la revisión III; originales preservados.
% Propietarios: U016, ecuaciones extractor-diferencias-k y lector-armonico-*;
% registro_k.tex, subsecciones k-pi-h, k-hadamard y k-transversal.
\subsection{Compatibilidad integral y determinación del registro transversal}
\label{img09:seccion}

El dominio de esta construcción es la publicación del registro de una
historia APP--TRIT--TPK. Se distinguen la evaluación de los eventos del
registro, la compatibilidad de sus canales y la inversión de sus
coordenadas. El resultado siguiente da un criterio integral completo para
el segundo nivel y una factorización explícita para conectar el primero con él.
La demostración utiliza los operadores del registro y establece sus
relaciones para todo elemento de los dominios enteros indicados.

\subsubsection{Caracterización por incrementos adyacentes}

Los índices pertenecen a $\mathbb Z/12\mathbb Z$. Sean
\[
 (Sz)_i=z_{i+1},\qquad D_3=I-S^3,\qquad D_4=I-S^4,
 \qquad Q(z)=\sum_{i=0}^{11}z_i.
\]
Se conserva el extractor transversal
\[
 \mathcal D:\mathbb Z^{12}\longrightarrow\mathbb Z^{25},
 \qquad k\longmapsto(D_3k,D_4k,Q(k)).
\]
Para una terna de coordenadas enteras $(b,c,q)$, escribimos
\begin{equation}
 w_i=c_i-b_{i+1},\qquad
 P_0=0,\quad P_i=\sum_{j=0}^{i-1}w_j\ (1\leq i\leq12),
 \qquad T(w)=\sum_{i=0}^{11}P_i.
 \label{img09:incremento}
\end{equation}

\begin{theorem}[Imagen integral del extractor transversal]
\label{img09:imagen}
Una terna $(b,c,q)\in\mathbb Z^{25}$ pertenece a
$\operatorname{im}\mathcal D$ si y sólo si satisface
\begin{align}
 b_i&=w_i+w_{i+1}+w_{i+2}&& (0\leq i<12),
 \label{img09:relaciones}\\
 \sum_{i=0}^{11}w_i&=0,
 \label{img09:cierre}\\
 q+T(w)&\equiv0\pmod {12}.
 \label{img09:congruencia}
\end{align}
Su única preimagen entera es
\begin{equation}
 k_i=\frac{q+T(w)}{12}-P_i,\qquad0\leq i<12.
 \label{img09:inversa}
\end{equation}
Las doce ecuaciones de \eqref{img09:relaciones} y la ecuación
\eqref{img09:cierre} son linealmente independientes sobre $\mathbb Q$.
\end{theorem}
\begin{proof}
Para una salida de $\mathcal D$ se tiene
\[
 c_i-b_{i+1}=(k_i-k_{i+4})-(k_{i+1}-k_{i+4})=k_i-k_{i+1}.
\]
La suma de tres incrementos da $b_i$; la suma cíclica de todos ellos
es cero. Además, $P_i=k_0-k_i$, por lo que
$q=12k_0-T(w)$. Se obtienen todas las condiciones y la fórmula inversa.

Recíprocamente, \eqref{img09:congruencia} hace entera la fórmula
\eqref{img09:inversa}; \eqref{img09:cierre} permite prolongarla
cíclicamente. Sus diferencias adyacentes son $w_i$. Por tanto,
$D_3k=b$ y
\[
 (D_4k)_i=w_i+w_{i+1}+w_{i+2}+w_{i+3}
 =w_i+b_{i+1}=c_i.
\]
La suma de la inversa da $Q(k)=q$. La misma fórmula prueba la unicidad.
Para la independencia, el cambio entero invertible
$(b,c)\mapsto(b,w=c-Sb)$ convierte las doce primeras relaciones en
$b=(I+S+S^2)w$. Cada una despeja una coordenada diferente de $b$.
La condición $\sum w_i=0$ es una relación adicional independiente.
\end{proof}

La imagen racional tiene dimensión doce. Su estructura integral añade
una congruencia de orden doce: ésta es la manifestación constructiva del
último factor de Smith de $\mathcal D$. Los once incrementos
$w_0,\ldots,w_{10}$ y la carga $q$ son coordenadas suficientes, con
$w_{11}=-\sum_{i=0}^{10}w_i$ y la congruencia
\eqref{img09:congruencia}. Las veinticinco coordenadas originales
conservan esos mismos doce grados de libertad mediante trece relaciones.

\subsubsection{Conmutatividad con la lectura armónica}

Sea $\mathcal H_{12}$ la transformación de Hadamard del registro,
con órbitas $(r,r+3,r+6,r+9)$ para $r=0,1,2$. Su bloque es
\[
 H_4=\begin{pmatrix}
 1&1&1&1\\1&1&-1&-1\\1&-1&1&-1\\1&-1&-1&1
 \end{pmatrix},\qquad H_4^2=4I_4.
\]
Para $(b,c,q)$ compatible, se definen
\[
 B_r=\sum_{j=0}^3c_{r+3j},\qquad
 A_0=\frac{q+2B_0+B_1}{3},\quad A_1=A_0-B_0,
 \quad A_2=A_1-B_1.
\]
El lector $\Pi_H$ de \eqref{eq:k-bloques-firmados} tiene los bloques
\[
 (\Pi_H(b,c,q))^{(r)}=
 (A_r,b_{r+3}-b_{r+9},b_r+b_{r+6},b_r-b_{r+6}).
\]

\begin{proposition}[Concordancia de las dos reconstrucciones]
\label{img09:triangulo}
Sobre $\operatorname{im}\mathcal D$ se cumple
\[
 \Pi_H\mathcal D=\mathcal H_{12},\qquad
 \mathcal D^{-1}=\tfrac14\mathcal H_{12}\Pi_H.
\]
La imagen de $\Pi_H$ restringida a ese dominio es exactamente
\[
 \mathcal L_H=
 \{u\in\mathbb Z^{12}:\mathcal H_{12}u\equiv0\pmod4\}.
\]
\end{proposition}
\begin{proof}
Para $k=\mathcal D^{-1}(b,c,q)$, sea
$a_r=\sum_{j=0}^3k_{r+3j}$. El desplazamiento de cuatro posiciones
lleva cada órbita a la siguiente, y por ello $B_r=a_r-a_{r+1}$.
Junto con $q=a_0+a_1+a_2$, estas identidades dan $A_r=a_r$.
En una órbita, escribamos $(x_0,x_1,x_2,x_3)$ para las coordenadas
de $k$ y $d_j=x_j-x_{j+1}$. Entonces
\[
 d_1-d_3=x_0+x_1-x_2-x_3,\quad
 d_0+d_2=x_0-x_1+x_2-x_3,\quad
 d_0-d_2=x_0-x_1-x_2+x_3.
\]
Son las tres filas restantes de $H_4k^{(r)}$. La inversión y la
caracterización integral se siguen de $\mathcal H_{12}^2=4I$.
\end{proof}

La condición $\Pi_H(b,c,q)\in\mathcal L_H$ sólo controla la
salida armónica. La pertenencia de las veinticinco coordenadas a
$\operatorname{im}\mathcal D$ exige además
\eqref{img09:relaciones}--\eqref{img09:congruencia}.
Por ejemplo, añadir a $c$ el vector $e_0-e_3$ conserva las tres sumas
orbitales y deja $\Pi_H$ inalterado. Partiendo de $(0,0,0)$ se
obtiene así una terna con $\Pi_H=0$ que incumple
\eqref{img09:relaciones}. Éste es un control negativo exacto de la
compatibilidad del canal, independiente de toda evaluación terminal.

\subsubsection{Determinación y transporte sobre las órbitas del registro}

El Teorema~\ref{img09:imagen} determina $k$ cuando se han producido
sus canales y su carga. Como condición sobre salidas aún no
individualizadas, la misma imagen es un grupo isomorfo a
$\mathbb Z^{12}$. Fijada únicamente una carga $q$, las soluciones
forman una clase afín de
\[
 L_0=\{v\in\mathbb Z^{12}:\sum v_i=0\}
 =\bigoplus_{i=0}^{10}\mathbb Z(e_i-e_{11}).
\]
Por ejemplo, $100\mathbf1+e_0$ y $100\mathbf1+e_1$ son dos
registros distintos, de carga $1201$, con componentes en
$\{0,\ldots,999\}$. Cada uno satisface todas las condiciones de
imagen y todas las congruencias de Hadamard mediante sus propios
canales. Estos testigos comparan las ecuaciones de compatibilidad;
no se les atribuye la condición de historias terminales HMT.

\begin{proposition}[Covariancia cíclica y estabilizador del registro]
\label{img09:covariancia}
Sean $\rho$ la traslación de ventanas en un dominio de registros
$\mathscr R$ y $S$ la traslación anterior en $\mathbb Z^{12}$.
Sobre la órbita de $R\in\mathscr R$, un lector covariante queda
determinado por un vector
\[
 k\in\operatorname{Fix}(S^p),
\]
donde $p\mid12$ es el período mínimo de $R$ bajo $\rho$.
La condición de carga $Q(k)=q$ admite soluciones enteras si y sólo si
$12/p$ divide a $q$; cuando existen, forman una clase afín de rango
$p-1$. En particular, sobre una órbita libre, covariancia y carga
dejan once grados de libertad.
\end{proposition}
\begin{proof}
La regla $F(\rho^jR)=S^jk$ está bien definida exactamente cuando
$S^pk=k$. Esos vectores repiten $12/p$ veces una palabra de
longitud $p$. Su carga es $\frac{12}{p}\sum_{i=0}^{p-1}k_i$.
La divisibilidad y el rango se siguen de esta fórmula. La
covariancia de los canales resulta de $D_aS=SD_a$ y $QS=Q$.
Sobre $\mathcal L_H$, la acción correspondiente es
$\widehat S_H=\mathcal H_{12}S\mathcal H_{12}/4$;
por la Proposición~\ref{img09:triangulo}, ambas reconstrucciones
entrelazan estas acciones.
\end{proof}

Si el dominio lleva un origen marcado, debe utilizarse la acción que
transporta también esa marca. La proposición se refiere a esa acción
cíclica declarada; no presupone la periodicidad de la historia
enriquecida. La naturalidad frente a truncamientos posteriores al corte
$108$ se obtiene, para cualquier lector $F_{108}$ ya definido, mediante
$F_N=F_{108}\circ\tau_{N,108}$. Esta naturalidad conserva la regla
de evaluación inicial; por sí sola no elige una de las reglas posibles
sobre los primeros ciento ocho eventos.

\subsubsection{Evaluación por eventos y composición de contribuciones}

La caracterización anterior proporciona una factorización de la extracción mediante dos
lecturas efectivas del registro:
\begin{equation}
 \omega:\mathscr R\longrightarrow
 \{w\in\mathbb Z^{12}:\sum w_i=0\},\qquad
 q_{\rm reg}:\mathscr R\longrightarrow\mathbb Z,
 \quad q_{\rm reg}(R)+T(\omega(R))\equiv0\pmod {12}.
 \label{img09:contrato-productor}
\end{equation}
Estas aplicaciones deben evaluarse sobre los eventos, la orientación,
la incidencia y la memoria que produjo APP--TRIT--TPK. Una vez dadas sus
acciones sobre esos datos, la composición
\[
 R\longmapsto(w,q)\longmapsto
 ((I+S+S^2)w,(I+S+S^2+S^3)w,q)
 \xrightarrow{\Pi_H}u\xrightarrow{\mathcal H_{12}/4}k
\]
está determinada y satisface todos los controles inversos. La fórmula
$w=c-Sb$ también permite conectar un productor que ya entregue ambos
canales: es una comprobación sobre su salida, anterior a toda comparación
con el testigo terminal.

\begin{proposition}[Acumulación de contribuciones y unicidad prospectiva]
\label{img09:acumulacion}
Supóngase que el estado inicial y cada evento $a_t$ del registro
producen, mediante reglas fijadas previamente, contribuciones
$(\delta w_t,\delta q_t)$ con
\[
 \sum_i\delta w_{t,i}=0,\qquad
 \delta q_t+T(\delta w_t)\equiv0\pmod {12}.
\]
Con condiciones iniciales compatibles $(w^{(0)},q^{(0)})$, la
actualización
\[
 w^{(t+1)}=w^{(t)}+\delta w_t,\qquad
 q^{(t+1)}=q^{(t)}+\delta q_t
\]
produce un único registro en cada etapa. Su incremento es
\[
 \delta k_{t,i}=
 \frac{\delta q_t+T(\delta w_t)}{12}-P_i(\delta w_t).
\]
La construcción respeta la concatenación de registros, y la lectura
de un prefijo permanece igual bajo prolongaciones posteriores.
\end{proposition}
\begin{proof}
Las condiciones de compatibilidad se conservan por suma. La fórmula
\eqref{img09:inversa} es aditiva en $(w,q)$ y produce el
incremento indicado. La inducción fija sucesivamente cada estado.
La suma sobre dos segmentos es la suma concatenada, y un prefijo recibe
exactamente los mismos sumandos antes y después de prolongar la historia.
\end{proof}

La proposición proporciona un criterio suficiente, no una exigencia de
linealidad para todo lector TPK. Si los eventos se transportan por acciones
no triviales, sus contribuciones pueden expresarse en la carta terminal
mediante el cociclo afín correspondiente. La ecuación de cociclo conserva
entonces la misma independencia de la partición de la trayectoria.
En ambos casos, las reglas de evento y el estado inicial constituyen los
datos constructivos que han de provenir del artículo: escogerlos para
reproducir un vector esperado sustituiría esa procedencia por otra tarea.

El resultado de este desarrollo es la condición exacta de determinación
\eqref{img09:contrato-productor}, su inversión demostrada y sus identidades
de transporte. Su alcance es el registro transversal; las restantes
construcciones de HMT conservan sus dominios y sus pruebas propias.
